%%%%%%%%%%%%%%%%%%%%%%%%%%%%%%%%%%%%%%%%%%%%%%%%%%%%%%%%%%%%%%%%%%%%%%
% Lớp 11 | Hàm số lượng giác và phương trình lượng giác | Phương trình lượng giác cơ bản
% Dạng 1D1?5-1: Câu hỏi lý thuyết. Khái niệm phương trình tương đương
% Số câu: 8  (sinh tự động bởi tach_cau.py; từ nay thầy sửa trực tiếp file này)
%%%%%%%%%%%%%%%%%%%%%%%%%%%%%%%%%%%%%%%%%%%%%%%%%%%%%%%%%%%%%%%%%%%%%%

%%%%%%%%%%%%%%%% TRẮC NGHIỆM 4 PHƯƠNG ÁN %%%%%%%%%%%%%%%%

% nguồn: Ok_Toan11_TracNghiem_Lan2
\begin{ex}%[1D1N5-1]%[ma:91b37626]
	Trong các phương trình sau, phương trình  nào có nghiệm?
	\choice 
		{$\cos x=3$} 
		 {$\sin 2 x=-2$}
		{\True  $\cos \left(2 x-\dfrac{\pi}{3}\right)=-1$}
		{$\cos (2 x-1)=\dfrac{\sqrt{7}}{2}$} 
	 \loigiai{
		Phương trình lượng giác cơ bản dạng $\sin u=\alpha$, $\cos u=\alpha$ có nghiệm khi và chỉ khi $|\alpha| \leq 1$. Nên phương trình $\cos \left(2  x-\dfrac{\pi}{3}\right)=-1$ có nghiệm.
	} 
\end{ex}

% nguồn: Ok_Toan11_TracNghiem_Lan2
\begin{ex}%[1D1N5-1]%[ma:1d0d21db]
	Phương  trình nào sau đây có nghiệm?
	\choice 
		 {$\sin 2021x-2=0$}
		 {$\cos (2 x+2021)=3$}
		 {$\sin^{2}x+1=0$}
		{\True  $\cos(2x+2021)=-1$}
	 \loigiai{
		Phương trình lượng giác cơ bản dạng $\sin u=\alpha$, $\cos u=\alpha$ có nghiệm khi và chỉ  khi $|\alpha| \leq 1$. Nên phương trình $\cos(2x+2021)=-1$ có nghiệm.\\
		Do $\sin^{2}x+1 \ge 1 > 0$ nên phương trình $\sin^{2}x+1=0$ vô  nghiệm.
	} 
\end{ex}

% nguồn: Ok_TracNghiemLop11_Lan1
\begin{ex}%[1D1N5-1]%[ma:fbd89e92]
	Nghiệm của phương trình lượng giác $\sin  x=5$ là
	\choice 
		{$x\in \mathbb{R}$} 
		{$x=\pm \arcsin 5+k2\pi \left( k\in  \mathbb{Z} \right)$}
		{$\left[\begin{aligned}
		 & x=\arcsin 5+k2\pi \\
		& x=\pi -\arcsin 5+k2\pi \\ 
\end{aligned} \right.\left( k\in \mathbb{Z} \right)$}
 {\True $x\in \varnothing $}
 \loigiai{
	Phương trình $\sin x=m$ chỉ có nghiệm khi và chỉ khi $\left|  m \right|\le 1$.\\
	Nên phương trình đã cho vô nghiệm. 
 }
\end{ex}

% nguồn: Ok_TracNghiemLop11_Lan1
\begin{ex}%[1D1N5-1]%[ma:5a815c9e]
	Phương trình $\sin x=\sin \beta$ có nghiệm là 
	 \choice
		{$\hoac{&x=\beta+k \pi \\ 
		&x=-\beta+k \pi}; k \in \mathbb{Z}$} 
		{$\hoac{&x=\beta+k 2\pi \\ 
		&x=-\beta+k 2\pi}; k  \in \mathbb{Z}$}
		{$\hoac{&x=\beta+k \pi  \\
		&x=\pi-\beta+k \pi};  k \in \mathbb{Z}$}
		{\True $\hoac{&x=\beta+k 2\pi  \\
		&x=\pi-\beta+k 2\pi}; k \in \mathbb{Z}$} 
	\loigiai{ 
		 Ta có $\sin x=\sin \beta \Leftrightarrow \hoac{&x=\beta+k 2\pi \\
		&x=\pi-\beta+k 2\pi}; k  \in \mathbb{Z}$
	 }
\end{ex}

% nguồn: Ok_Toan11_TracNghiem_Lan2
\begin{ex}%[1D1V5-1]%[ma:2388b78b]
	Phương trình $  8\sin^2\left(\dfrac{x}{2}\right){\cos^2}\left(\dfrac{x}{2}\right)-1=0$ tương đương với phương trình nào sau đây?
	\choice 
		 {$\sin x=\dfrac{\sqrt{2}}{2}$}
		 {\True $\cos 2x=0$}
		{$\cos x=\dfrac{\sqrt{2}}{2}$} 
		{$\sin  x=-\dfrac{\sqrt{2}}{2}$}
	\loigiai{ 
		Ta  có
		\begin{eqnarray*}
			8\sin^2\left(\dfrac{x}{2}\right){\cos^2}\left(\dfrac{x}{2}\right)-1=0&\Leftrightarrow&  2\cdot \left(2\sin\left(\dfrac{x}{2}\right)\cos\left(\dfrac{x}{2}\right)\right)^2-1=0\Leftrightarrow 2\sin^2x-1=0\\
			&\Leftrightarrow& 1-2\sin^2x=0\Leftrightarrow\cos  2x=0.
		\end{eqnarray*}
		Vậy phương trình đã  cho tương đương với phương trình $\cos 2x=0$.
	 }
\end{ex}


%%%%%%%%%%%%%%%% ĐÚNG/SAI %%%%%%%%%%%%%%%%

% nguồn: Ok_dungsai_lop11_Lan1
\begin{ex}%[1D1N5-1]%[ma:6248ef2f]
 Phép biến đổi nào sau đây là  phép biến đổi tương đương?
 \choiceTF 
  {\True  $x+\dfrac{1}{x}=1+\dfrac{1}{x}\Leftrightarrow x=1$}
  {$\sqrt{x+2}=x-4\Leftrightarrow x+2=(x-4)^2$} 
  {$x(x+2)=x\Leftrightarrow x+2=1$} 
  {\True  $x+\sqrt{x-1}=x^2\Leftrightarrow x=x^2-\sqrt{x-1}$}
 \loigiai{\begin{itemize}
  \item Phương trình $x+\dfrac{1}{x}=1+\dfrac{1}{x}$ xác định khi $x\ne0$, do đó $x+\dfrac{1}{x}=1+\dfrac{1}{x}\Leftrightarrow  x=1$.
  \item $\sqrt{x+2}=x-4\Leftrightarrow x+2=(x-4)^2$ không là phép biến đổi tương đương, vì làm sinh ra nghiệm  ngoại lai $x=2$.
  \item $x(x+2)=x\Leftrightarrow x+2=1$ không là phép  biến đổi tương đương vì làm mất nghiệm $x=0$.
  \item $x+\sqrt{x-1}=x^2\Leftrightarrow x=x^2-\sqrt{x-1}$ là phép biến đổi tương đương  vì không làm ảnh hưởng đến tập nghiệm của bài toán.
 \end{itemize}
 }
\end{ex}

% nguồn: Ok_dungsai_lop11_Lan1
\begin{ex}%[1D1N5-1]%[ma:19d63c40]
 Cho hàm số $y=2\sin x\cos x$. 
 \choiceTF 
  {Rút gọn  hàm số đã cho, ta được $y=\cos 2x$}
  {Hàm số đã cho có tập tập giá trị là $\mathbb{R}$} 
  {\True Khi hàm số $y=0$ hay phương trình $2\sin x\cos x=0$ thì $x=0$ là một nghiệm của  phương trình}
  {Khi hàm số $y=1$ hay phương trình $2\sin x\cos x=1$ thì $x=2\pi$ là  một nghiệm của phương trình}
 \loigiai{ 
  \begin{enumerate}
   \item \textbf{Sai.} 
   Vì $y=2\sin  x\cos x=\sin2x$.
   \item \textbf{Sai.} 
   Vì tập  giá trị là $[-1;1]$ hoặc tập xác định là $\mathbb{R}$.
    \item \textbf{Đúng.}
   Vì thay $x=0$ vào phương  trình, ta được $y=2\sin0\cos0\Rightarrow y=0$.
    \item \textbf{Sai.}
   Vì thay $x=2\pi$ vào  phương trình, ta được $y=2\sin2\pi\cos2\pi=0\Rightarrow y\neq1$.
  \end{enumerate}
  }
\end{ex}

% nguồn: Ok_dungsai_lop11_Lan1
\begin{ex}%[1D1H5-1]%[ma:fef67c10]
 Trong các phương trình sau, phương trình nào tương đương với  phương trình $x^2=9$?
 \choiceTF 
   {$x-3=0$}
  {\True  $x^2+\sqrt{x^2}=9+\sqrt{x^2}$}
   {$x^3-9x=0$}
  {\True  $|x|=3$}
 \loigiai{Phương trình  $x^2=9$ có tập nghiệm là $S=\{-3;3\}$.
  \begin{itemize}
   \item Phương trình $x-3=0$ có tập  nghiệm là $S=\{3\}$ nên không tương đương với phương trình $x^2=9$.
   \item Do $x^2\ge0$, $\forall x\in\mathbb{R}$ nên phương trình $x^2+\sqrt{x^2}=9+\sqrt{x^2}$ tương đương  với $x^2=9$.
   \item Phương trình $x^3-9x=0$ có tập nghiệm là $S=\{-3;0;3\}$ nên không tương đương với phương trình $x^2=9$. 
   \item Phương trình $|x|=3$  có tập nghiệm là $S=\{-3;3\}$ nên tương đương với phương trình $x^2=9$.
 \end{itemize}}
\end{ex}
